\documentclass[10pt,letterpaper]{article}
\usepackage{cloudsuite}
\renewcommand{\DocShort}{ENGINEERING GUIDE}
\begin{document}
\coverpage{Cloud Computing Engineering Guide}{Vendor-neutral foundations, mechanisms and design decisions}{Concepts \textbullet{} decisions \textbullet{} verification artifacts}
\sourceStatement
\tableofcontents
\newpage
\section{Executive overview}
A cloud solution is a negotiated combination of on-demand resources, shared operational capabilities, contracts and responsibilities. Engineering rigor comes from tracing each claimed benefit to a capability, each capability to an operational owner and each risk to verification evidence. This guide follows the text's progression from fundamentals to mechanisms, architectures and service obligations.
\begin{takeaway}
\textbf{Design equation (conceptual):} Requirement $\rightarrow$ delivery model $\rightarrow$ mechanisms $\rightarrow$ architecture $\rightarrow$ control owner $\rightarrow$ measured service outcome. A cloud product name is neither a requirement nor an acceptance test.\end{takeaway}
\section{System context and actor model}
The book distinguishes cloud provider, cloud consumer, broker, service owner and resource administrator. Assign each role even when one legal entity holds several. Draw organizational boundaries separately from trust boundaries; a single organization can administer multiple trust zones.\\
\sourceRef{4}{51}
\begin{center}
\begin{tikzpicture}[node distance=0.55cm and 0.65cm, font=\footnotesize,
 box/.style={draw=navy,fill=pale,rounded corners,minimum height=1.15cm,text width=3cm,align=center},
 arr/.style={-Latex,thick,draw=petrol}]
 \node[box] (user) {Cloud Consumer\\Requirements and data ownership};
 \node[box,right=of user] (service) {Cloud Service\\Contract and workload};
 \node[box,right=of service] (provider) {Cloud Provider\\Platform and operations};
 \draw[arr] (user) -- (service);
 \draw[arr] (service) -- (provider);
 \node[draw=gold,dashed,fit=(service),inner sep=7pt,label={[font=\scriptsize,color=gold]below:Trust and responsibility boundary}] {};
\end{tikzpicture}
\end{center}
\subsection{Minimum context record}
\begin{itemize}
\item Consumer and service owner; provider and broker if present; resource administrators and privileged operators.
\item Application data classes, locations and data transfer boundaries.
\item Chosen delivery model and deployment model; identity federation and network crossings.
\item SLO/SLA target, commercial metering unit and recovery expectations.
\end{itemize}
\section{Cloud characteristics and deployment choices}
\subsection{On-demand usage}
Provision resources through defined requests without waiting for a bespoke physical build. Operational evidence should include an owner, a scope and a repeatable test.
\subsection{Ubiquitous access}
Access resources through appropriate network-facing interfaces with governed identity. Operational evidence should include an owner, a scope and a repeatable test.
\subsection{Multitenancy and pooling}
Allocate shared capacity while preventing cross-consumer disclosure and unacceptable contention. Operational evidence should include an owner, a scope and a repeatable test.
\subsection{Elasticity}
Match capacity to changing demand under limits and policy. Operational evidence should include an owner, a scope and a repeatable test.
\subsection{Measured usage}
Record consumption in a form suitable for allocation, forecasting or charging. Operational evidence should include an owner, a scope and a repeatable test.
\subsection{Resiliency}
Sustain or restore service under specified failures and degraded conditions. Operational evidence should include an owner, a scope and a repeatable test.
\subsection{Delivery-model responsibility matrix}
\small\begin{longtable}{p{0.2\linewidth}p{0.23\linewidth}p{0.23\linewidth}p{0.23\linewidth}}
\toprule Control plane or asset & IaaS & PaaS & SaaS \\\midrule
Physical hosts / hypervisor & Provider & Provider & Provider \\
Guest OS or managed runtime & Consumer manages guest OS & Provider manages platform runtime & Provider manages application runtime \\
Application logic & Consumer & Consumer & Provider \\
Data classification and lawful use & Consumer accountability & Consumer accountability & Consumer accountability \\
Application access configuration & Consumer & Consumer & Consumer (available controls) \\
Audit and evidence & Contract-specific split & Contract-specific split & Contract-specific split \\\bottomrule
\end{longtable}\normalsize
A responsibility matrix states the usual direction of operational control; specific contracts and configuration choices take precedence. Neither PaaS nor SaaS removes consumer accountability for identity governance or data handling.
\section{Network, virtualization and container foundations}
\subsection{Network fabric}
Model the client path, routing, WAN or internet dependency, bandwidth, packet loss and latency budget; validate failure at the boundary.
\subsection{Virtual servers}
Specify host/guest patch responsibility, hypervisor isolation, image baseline, disk attachment and remote access.
\subsection{Containers}
Control build provenance, image digests, software inventory, runtime permissions, pod network policies and orchestration behavior.
\subsection{Service interfaces}
Treat REST or RPC APIs as contracts: authentication, authorization, schemas, timeouts, idempotency, versions and errors.
\section{Mechanism composition patterns}
\subsection{Elastic stateless service}
Use a cloud usage monitor to observe demand, an automated scaling listener to interpret thresholds, resource replication to change instance count and a load balancer to route only healthy requests. Check latency while adding capacity and during scale-in.
\subsection{Recoverable stateful service}
Combine storage, data backup and recovery, activity monitoring, failover and a state management database where the service requires durable coordination. Verify write durability, replay semantics, data loss and restart outcomes.
\subsection{Auditable privileged administration}
Connect IAM, MFA, remote administration, audit and authentication logging. Deny unattended privileged pathways and require a retrievable session trail.
\subsection{Metered managed service}
Correlate cloud usage, pay-per-use metering, service-quality monitoring, billing management and SLA management. Define the source of truth for each reported value.
\section{Security model and threat-to-control mapping}
The threat taxonomy includes eavesdropping, malicious intermediaries, denial of service, insufficient authorization, attacks against virtualization or containers, malware, privilege escalation and code/data attacks. Distinguish preventative controls from detective and recovery mechanisms.\\
\sourceRef{7}{159}
\small\begin{longtable}{p{0.22\linewidth}p{0.3\linewidth}p{0.36\linewidth}}
\toprule Threat scenario & Control composition & Required evidence \\\midrule
Stolen operator credentials & IAM; MFA; SSO; authentication log monitor & Revocation test, log events, privilege inventory \\
Unapproved network path & Logical network perimeter; firewall; security groups & Explicit allow/deny traffic test \\
Malicious or outdated workload & Hardened image; update utility; malicious code analysis & Artifact provenance, compliance and scan records \\
Sensitive data movement & Encryption; DLP; data loss monitor & Boundary classification and controlled egress test \\
Availability disruption & Failover; load balancer; backup and recovery & Measured disruption, switch-over and recovery \\
Audit blind spot & Audit monitor; activity log monitor & Completeness, retention and tamper detection \\\bottomrule
\end{longtable}\normalsize
\section{Selecting cloud architectures}
The source introduces 38 named architectures in three groups. Start with an explicit system quality, then consider a fundamental architecture; adopt advanced or specialized arrangements only when their additional operational assumptions are justified.
\subsection{Fundamental selection questions}
\begin{itemize}
\item Which requirement cannot be met by a simpler arrangement?
\item Which prerequisite mechanisms and failure domains will this introduce?
\item Which operator can test and maintain it, and how often?
\item What measurable acceptance criterion will prove value?
\end{itemize}
\emph{Representative source concepts:} Workload Distribution, Resource Pooling, Dynamic Scalability, Elastic Resource Capacity, Service Load Balancing, Cloud Bursting, Elastic Disk Provisioning, Redundant Storage, Multicloud.\\
\sourceRef{13}{343}
\subsection{Advanced selection questions}
\begin{itemize}
\item Which requirement cannot be met by a simpler arrangement?
\item Which prerequisite mechanisms and failure domains will this introduce?
\item Which operator can test and maintain it, and how often?
\item What measurable acceptance criterion will prove value?
\end{itemize}
\emph{Representative source concepts:} Hypervisor Clustering, Virtual Server Clustering, Zero Downtime, Resilient Disaster Recovery, Distributed Data Sovereignty, Rapid Provisioning, Virtual Private Cloud.\\
\sourceRef{14}{371}
\subsection{Specialized selection questions}
\begin{itemize}
\item Which requirement cannot be met by a simpler arrangement?
\item Which prerequisite mechanisms and failure domains will this introduce?
\item Which operator can test and maintain it, and how often?
\item What measurable acceptance criterion will prove value?
\end{itemize}
\emph{Representative source concepts:} Direct I/O Access, Storage tiering, Multipath Resource Access, Edge Computing, Fog Computing, Virtual Data Abstraction, Federated Cloud Application.\\
\sourceRef{15}{415}
\section{Cloud provider and consumer operating model}
\begin{itemize}
\item Record which party provisions compute, storage and connectivity; who owns hypervisor, runtime, application and configuration updates.
\item Document access, monitoring and contractual escalation across the provider-consumer interface.
\item Make exit and portability explicit: data export, portability constraints, dependency removal and verification of residual access.
\item Require a deployment-model boundary review whenever an integration introduces a different provider, geography or organization.
\end{itemize}
\section{Cost engineering}
Use a workload-normalized model rather than a single hourly price. Under documented assumptions, the planning expression is:
\[
 C_{\mathrm{period}}=C_{\mathrm{one\mbox{-}time}}+\sum_i (q_i\,r_i)
 +C_{\mathrm{operations}}+C_{\mathrm{transfer}}+C_{\mathrm{risk}}.
\]
Here $q_i$ is metered quantity, $r_i$ its unit price, and the final terms represent explicitly estimated operational, data-transfer and risk allowances. This is an instructional model, not a vendor invoice formula. Assess on-demand allocation, reserved commitments, outbound transfer, storage and service subscriptions separately.\\
\sourceRef{17}{479}
\section{Availability, recovery and SLAs}
\[
 A=1-\frac{T_{\mathrm{eligible\ outage}}}{T_{\mathrm{eligible\ observation}}},\qquad
 \mathrm{RPO}=\text{maximum acceptable data-loss age},\quad
 \mathrm{RTO}=\text{target restoration duration}.
\]
Define measurement eligibility, exclusions and evidence owner before reporting a service-level percentage. Distinguish failure frequency, mean time between failures, response latency, scaling capacity and time to switchover; one metric cannot substitute for all others. RPO/RTO are common planning objectives used here as added operational context.\\
\sourceRef{18}{503}
\section{Architecture acceptance checklist}
\begin{itemize}
\item Requirements and rationale trace to the selected delivery and deployment models.
\item Actors, organizational boundaries and trust boundaries are explicitly diagrammed.
\item Infrastructure, specialized, security and management mechanisms are owned and verified.
\item Availability/recovery and usage/cost targets have unambiguous measurement definitions.
\item Failure injection, identity denial, data restoration and cost attribution tests pass.
\item Operations, contractual escalation, exception handling and exit paths are documented.
\end{itemize}
\begin{workbox}[title={Integrated capstone}]
Design an API-based service with personal data, seasonal traffic and a contractual uptime target. Produce: a context diagram, a service-model RACI, one scale-out path, one control-to-threat map, one recovery test, one metered cost scenario and an evidence-based acceptance checklist.
\end{workbox}

\refs
\end{document}
